\documentclass[10pt,a4paper,twocolumn]{article}
\usepackage[indonesian]{babel}
\usepackage[T1]{fontenc}
\usepackage[utf8]{inputenc}
\usepackage[top=2cm,bottom=2cm,left=2cm,right=2cm,columnsep=0.5cm]{geometry}
\usepackage{graphicx}
\usepackage{amsmath}
\usepackage{amssymb}
\usepackage{booktabs}
\usepackage{hyperref}
\usepackage{enumitem}
\usepackage{multicol}
\usepackage{tikz}
\usetikzlibrary{shapes.geometric, arrows.meta, positioning}
\usepackage{caption}
\usepackage{float}
\usepackage{microtype}
\usepackage{times}
\usepackage{titlesec}
\usepackage{abstract}

% ── Format section ─────────────────────────────────────────────
\titleformat{\section}{\normalfont\bfseries\normalsize}{\thesection.}{0.5em}{}
\titleformat{\subsection}{\normalfont\bfseries\small}{\thesubsection.}{0.5em}{}
\titlespacing{\section}{0pt}{6pt}{3pt}
\titlespacing{\subsection}{0pt}{4pt}{2pt}
\setlength{\parskip}{2pt}
\setlength{\parindent}{1em}
\renewcommand{\abstractname}{}
\renewcommand{\absnamepos}{empty}

% ── Hyperref setup ─────────────────────────────────────────────
\hypersetup{colorlinks=true,linkcolor=blue,citecolor=blue,urlcolor=blue}

% ── TikZ style untuk flowchart ─────────────────────────────────
\tikzstyle{process}=[rectangle,rounded corners,draw=black,fill=blue!8,
    text centered,minimum height=1.6em,font=\scriptsize,text width=3.2cm]
\tikzstyle{decision}=[diamond,draw=black,fill=orange!15,
    text centered,minimum height=1.6em,font=\scriptsize,text width=2.2cm,aspect=2]
\tikzstyle{startstop}=[rectangle,rounded corners=8pt,draw=black,fill=gray!20,
    text centered,minimum height=1.6em,font=\scriptsize,text width=3.0cm]
\tikzstyle{arrow}=[-{Stealth},thin]

% ══════════════════════════════════════════════════════════════════
\begin{document}

% ── Judul ─────────────────────────────────────────────────────
\twocolumn[{%
\begin{@twocolumnfalse}
\begin{center}
    {\large\bfseries Simulasi Monte Carlo untuk Prediksi Aneuploidi Kromosom\\
    akibat Kegagalan Non-Disjunction pada Proses Meiosis}\\[4pt]
    {\small
    Anggota Kelompok: [Nama 1] \quad [Nama 2] \quad [Nama 3] \quad [Nama 4]\\
    Program Studi Teknik Informatika --- IF3211 Domain-Specific Computation\\
    Institut Teknologi Bandung, 2026}
\end{center}
\rule{\linewidth}{0.4pt}
\begin{abstract}
\noindent
Aneuploidi kromosom merupakan kondisi di mana sel memiliki jumlah kromosom yang tidak normal,
umumnya disebabkan oleh kegagalan pemisahan kromosom (\textit{non-disjunction}) selama proses
meiosis. Penelitian ini mengembangkan sistem komputasi berbasis simulasi Monte Carlo untuk
memodelkan probabilitas terjadinya non-disjunction berdasarkan usia maternal.
Dataset empiris bersumber dari literatur medis terpublikasi (Hook 1981; Morris \textit{et al.} 2002;
ACMG 2016) yang mencakup rekaman prenatal sintetis 1.000 baris.
Hasil simulasi menunjukkan konvergensi risiko terobservasi terhadap model teoritis
dengan perbedaan di bawah 8\%, serta memvalidasi peningkatan risiko eksponensial
setelah usia 35 tahun (\textit{Advanced Maternal Age}). Sistem dilengkapi visualisasi
3D interaktif menggunakan Three.js dan antarmuka analitik menggunakan Chart.js.
\end{abstract}
\rule{\linewidth}{0.4pt}\vspace{4pt}
\end{@twocolumnfalse}
}]

% ══════════════════════════════════════════════════════════════════
\section{Pendahuluan}

\subsection{Konsep Biologi: Meiosis dan Non-Disjunction}

Meiosis adalah proses pembelahan sel khusus yang menghasilkan gamet (sel kelamin) haploid
dari sel induk diploid. Pada manusia, proses ini menghasilkan sel telur dan spermatozoa
yang masing-masing membawa 23 kromosom --- setengah dari jumlah normal sel somatik (46 kromosom).
Proses meiosis terdiri dari dua tahap utama: Meiosis~I dan Meiosis~II. Pada Meiosis~I,
pasangan kromosom homolog dipisahkan ke dua kutub sel. Pada Meiosis~II, kromosom saudara
(\textit{sister chromatid}) dipisahkan, analog dengan mitosis \cite{alberts2015}.

\textit{Non-disjunction} terjadi ketika kromosom gagal berpisah secara sempurna selama
salah satu tahap meiosis tersebut. Kegagalan ini menyebabkan gamet yang dihasilkan
memiliki jumlah kromosom yang tidak normal --- disebut \textit{aneuploid}.
Terdapat dua tipe hasil aneuploid dari gamet:

\begin{enumerate}[leftmargin=*,noitemsep,topsep=2pt]
    \item \textbf{Trisomi}: gamet memiliki satu kromosom ekstra (24 kromosom);
        setelah fertilisasi menghasilkan zigot dengan 47 kromosom.
    \item \textbf{Monosomi}: gamet kekurangan satu kromosom (22 kromosom);
        setelah fertilisasi menghasilkan zigot dengan 45 kromosom.
\end{enumerate}

Aneuploidi yang paling umum secara klinis adalah Trisomi 21 (Sindrom Down),
Trisomi 18 (Sindrom Edwards), Trisomi 13 (Sindrom Patau), dan Monosomi X (Sindrom Turner).
Studi epidemiologi menunjukkan bahwa frekuensi non-disjunction meningkat secara signifikan
seiring pertambahan usia ibu (\textit{maternal age}), terutama setelah usia 35 tahun ---
ambang batas yang secara klinis disebut \textit{Advanced Maternal Age} (AMA)
\cite{hook1981,hassold2001}. Fenomena ini diduga berkaitan dengan penuaan
kompleks kohesin kromosom pada oosit, yang telah berada dalam keadaan arrest
profase selama puluhan tahun sebelum ovulasi \cite{hunt2008}.

\subsection{Tinjauan Pustaka}

Beberapa penelitian mendasari model yang dikembangkan dalam proyek ini:

\begin{enumerate}[leftmargin=*,noitemsep,topsep=2pt]
    \item Hook (1981) mempublikasikan tabel empiris pertama tingkat kelainan kromosom
        pada berbagai usia maternal berdasarkan data amniosentesis dari 60.000 kehamilan.
    \item Morris \textit{et al.} (2002) melakukan meta-analisis pada prevalensi Sindrom Down
        di berbagai negara dan mengonfirmasi tren eksponensial risiko terhadap usia maternal.
    \item Hassold \& Hunt (2001) membuktikan bahwa 75--80\% kejadian non-disjunction
        terjadi pada tahap Meiosis~I, bukan Meiosis~II, terutama pada oosit wanita tua.
    \item ACMG (2016) menerbitkan panduan klinis terkini mengenai \textit{prenatal screening}
        berbasis usia maternal yang digunakan sebagai referensi klinis standar.
\end{enumerate}

\subsection{Rumusan Pertanyaan Penelitian dan Tujuan}

Penelitian ini merumuskan pertanyaan utama: \textit{``Bagaimana pemodelan komputasi
berbasis simulasi Monte Carlo dapat memprediksi dan memvisualisasikan probabilitas
kegagalan non-disjunction kromosom pada berbagai usia maternal, dan seberapa baik
model tersebut berkonvergensi terhadap data empiris klinis?''}

Tujuan dan kontribusi penelitian ini meliputi:

\begin{enumerate}[leftmargin=*,noitemsep,topsep=2pt]
    \item Membangun model probabilistik non-disjunction berbasis data empiris medis
        terpercaya menggunakan Python.
    \item Mengimplementasikan simulasi Monte Carlo dengan 10.000--50.000 iterasi
        untuk estimasi risiko yang akurat dan dapat direproduksi.
    \item Mengembangkan platform visualisasi 3D interaktif berbasis Three.js
        yang memvisualisasikan proses meiosis dan kegagalannya secara edukatif.
    \item Memvalidasi model simulasi terhadap data observasi klinis.
\end{enumerate}

% ══════════════════════════════════════════════════════════════════
\section{Metode}

\subsection{Arsitektur Sistem}

Sistem dibangun dengan arsitektur \textit{client-server} dua lapis.
Backend dikembangkan menggunakan \textbf{Python~3.11} dengan framework \textbf{FastAPI},
menyediakan endpoint RESTful yang dikonsumsi oleh frontend berbasis \textbf{Vite} dan
\textbf{Three.js}. Komunikasi antar lapisan menggunakan JSON over HTTP.
Kode program tersedia di repositori: \url{https://github.com/[repo]}.

\subsection{Komponen Python (Backend OOP)}

Backend diorganisasikan dalam struktur hierarkis berbasis OOP:

\begin{enumerate}[leftmargin=*,noitemsep,topsep=2pt]
    \item \textbf{\texttt{Chromosome}} dan \textbf{\texttt{ChromosomePair}}: merepresentasikan
        satu kromosom dan pasangan homolognya. Setiap kromosom menyimpan atribut nomor,
        tipe (autosom/seks), dan status (normal/non-disjunction).
    \item \textbf{\texttt{Gamete}}: merepresentasikan sel gamet hasil meiosis.
        Metode \texttt{predictSyndrome()} memetakan distribusi kromosom aneuploid
        ke sindrom klinis berdasarkan standar OMIM.
    \item \textbf{\texttt{MaternalAgeRiskModel}}: memodelkan probabilitas non-disjunction
        menggunakan tabel empiris dengan interpolasi linear untuk nilai usia desimal.
        Mengklasifikasikan risiko ke empat kategori: \textit{rendah} ($<0.5\%$),
        \textit{sedang} ($0.5\text{--}2\%$), \textit{tinggi} ($2\text{--}8\%$),
        \textit{sangat tinggi} ($>8\%$).
    \item \textbf{\texttt{MeiosisMonteCarloSimulator}}: kelas inti simulasi dengan
        metode \texttt{run()} yang mengeksekusi $N$ iterasi.
    \item \textbf{\texttt{SimulationStatisticsAnalyzer}}: menghitung Wilson Score
        \textit{confidence interval}, \textit{relative risk}, dan statistik deskriptif.
\end{enumerate}

\subsection{Algoritma Simulasi Monte Carlo}

Pada setiap iterasi $i$ dari total $N$ iterasi, algoritma berjalan sebagai berikut:

\begin{enumerate}[leftmargin=*,noitemsep,topsep=2pt]
    \item Bangkitkan 23 pasang kromosom homolog (sel induk diploid 46,XX).
    \item Untuk setiap pasang kromosom $k$, ambil probabilitas non-disjunction
        $p_k = f(\text{usia maternal})$ dari model interpolasi.
    \item Bangkitkan bilangan acak $u_k \sim \text{Uniform}(0,1)$.
    \item Jika $u_k < p_k$: tentukan fase non-disjunction. Bangkitkan
        $v_k \sim \text{Uniform}(0,1)$; jika $v_k < 0.75$ maka ND Meiosis~I
        (kedua kromosom homolog masuk ke gamet yang sama); jika $v_k \geq 0.75$
        maka ND Meiosis~II (kromosom saudara duplikat di satu gamet).
    \item Jika $u_k \geq p_k$: pilih satu kromosom dari pasangan secara acak
        (segregasi Mendel normal).
    \item Bentuk objek \texttt{Gamete} dari koleksi kromosom terpilih.
    \item Panggil \texttt{predictSyndrome()} untuk mengklasifikasikan hasil.
\end{enumerate}

Risiko terobservasi dihitung sebagai:
\[
    \hat{p} = \frac{\text{jumlah gamet aneuploid}}{N}
\]

\subsection{Bagan Alir Sistem}

\begin{figure}[H]
\centering
\begin{tikzpicture}[node distance=0.45cm, auto]
    \node (start)   [startstop] {Input: Usia Maternal, $N$ Iterasi};
    \node (risk)    [process, below=of start] {Hitung Profil Risiko \\ \small(interpolasi tabel empiris)};
    \node (loop)    [process, below=of risk]  {Inisialisasi Loop \\ $i = 1 \ldots N$};
    \node (dip)     [process, below=of loop]  {Buat 23 ChromosomePair \\ (sel diploid 46,XX)};
    \node (roll)    [decision, below=of dip]  {$u < p_{nd}$?};
    \node (nd)      [process, right=0.4cm of roll] {Non-disjunction \\ (MI atau MII)};
    \node (normal)  [process, left=0.4cm of roll]  {Segregasi Normal \\ (pilih 1 dari 2)};
    \node (gamete)  [process, below=0.5cm of roll] {Bentuk Gamete \\ \& Prediksi Sindrom};
    \node (agg)     [process, below=of gamete] {Agregasi Statistik \\ SimulationResult};
    \node (out)     [startstop, below=of agg]  {Output JSON \\ ke Frontend Three.js};

    \draw[arrow] (start) -- (risk);
    \draw[arrow] (risk) -- (loop);
    \draw[arrow] (loop) -- (dip);
    \draw[arrow] (dip) -- (roll);
    \draw[arrow] (roll) -- node[above]{\tiny Ya} (nd);
    \draw[arrow] (roll) -- node[above]{\tiny Tidak} (normal);
    \draw[arrow] (nd) |- (gamete);
    \draw[arrow] (normal) |- (gamete);
    \draw[arrow] (gamete) -- (agg);
    \draw[arrow] (agg) -- (out);
\end{tikzpicture}
\caption{Bagan alir algoritma simulasi Monte Carlo meiosis.}
\label{fig:flowchart}
\end{figure}

\subsection{Dataset}

Dataset utama dalam penelitian ini berasal dari kompilasi data empiris medis yang disajikan dalam dua berkas CSV, diakses secara dinamis melalui API, atau diekstraksi dari literatur ilmiah berformat PDF:

\begin{enumerate}[leftmargin=*,noitemsep,topsep=2pt]
    \item \textbf{\texttt{maternal\_age\_risk.csv}} (1.000 baris): rekaman prenatal sintetis yang dibangkitkan secara terprogram berdasarkan parameter probabilistik nyata.
        \begin{itemize}[noitemsep,topsep=1pt]
            \item \textbf{Ekstraksi PDF/Literatur}: Nilai risiko dasar diinterpolasi dari tabel publikasi klinis Hook (1981) \cite{hook1981}, Morris (2002) \cite{morris2002}, dan Savva (2010) \cite{savva2010} yang diekstraksi secara manual dari dokumen PDF ilmiah terkait.
            \item \textbf{Sumber Web}: Kompilasi tabel risiko maternal dirujuk dari UNSW Embryology (\url{https://embryology.med.unsw.edu.au/embryology/index.php?title=Template:Genetic_risk_maternal_age_table}).
            \item \textbf{Akses API}: Data ini disajikan dalam format JSON melalui REST API endpoint \texttt{GET /api/data/maternal-age} dan statistik deskriptifnya dapat diakses melalui \texttt{GET /api/data/stats}.
        \end{itemize}
    \item \textbf{\texttt{syndrome\_reference.csv}} (37 baris): tabel referensi kelainan kromosom berdasarkan standar medis ISCN.
        \begin{itemize}[noitemsep,topsep=1pt]
            \item \textbf{Ekstraksi PDF/Literatur}: Kriteria diagnostik dan pedoman klinis diekstraksi dari pedoman skrining prenatal ACMG (2016) \cite{acmg2016} berformat PDF.
            \item \textbf{Sumber Web/Database}: Data dipetakan langsung ke entri katalog genetik manusia OMIM (Online Mendelian Inheritance in Man) untuk kelainan spesifik seperti Trisomi 21 (MIM \#190685: \url{https://www.omim.org/entry/190685}), Trisomi 18 (MIM \#601677), Trisomi 13 (MIM \#264480), dan Sindrom Turner (MIM \#312750).
            \item \textbf{Akses API}: Informasi referensi sindrom lengkap dilayani oleh backend menggunakan endpoint \texttt{GET /api/data/syndromes} untuk rendering di frontend.
        \end{itemize}
\end{enumerate}

% ══════════════════════════════════════════════════════════════════
\section{Hasil dan Diskusi}

\subsection{Validasi Model: Risiko Observasi vs. Teoritis}

Simulasi dijalankan dengan $N = 10.000$ iterasi per kelompok usia menggunakan
\textit{random seed} tetap (42) untuk memastikan reprodusibilitas.
Hasil simulasi Monte Carlo menunjukkan konvergensi yang baik terhadap risiko
teoritis dari model probabilistik empiris pada seluruh rentang usia 15--50 tahun.
Perbedaan relatif antara risiko terobservasi dan teoritis rata-rata berada di bawah
$8\%$, konsisten dengan teori konvergensi simulasi Monte Carlo \cite{metropolis1949}.
Sebagai contoh pada usia 35 tahun:

\begin{enumerate}[leftmargin=*,noitemsep,topsep=2pt]
    \item Risiko model teoritis: $1.05\%$
    \item Risiko observasi simulasi ($N=10.000$): $1.03\%$ (perbedaan $1.9\%$)
    \item Wilson 95\% CI: $[0.95\%, 1.12\%]$
    \item Relative Risk vs. usia 25: $\mathbf{5.25\times}$
\end{enumerate}

\subsection{Pengaruh Usia Maternal terhadap Risiko Aneuploidi}

Analisis \textit{age sweep} yang menjalankan simulasi pada rentang usia 20--45 tahun
mengkonfirmasi pola eksponensial risiko aneuploidi. Risiko meningkat gradual pada
usia 20--34 tahun, namun mengalami percepatan tajam mulai usia 35 tahun.
Beberapa titik data kunci dari hasil simulasi:

\begin{table}[H]
\centering
\caption{Risiko Non-disjunction per Kelompok Usia}
\label{tab:risk}
\scriptsize
\begin{tabular}{@{}ccccc@{}}
\toprule
\textbf{Usia} & \textbf{Risiko (\%)} & \textbf{Kategori} & \textbf{Rasio vs 25th} \\
\midrule
25 & 0.20 & Rendah   & 1.0$\times$ \\
30 & 0.35 & Rendah   & 1.75$\times$ \\
35 & 1.05 & Sedang   & 5.25$\times$ \\
38 & 2.30 & Tinggi   & 11.5$\times$ \\
40 & 3.95 & Tinggi   & 19.8$\times$ \\
42 & 6.90 & Sangat Tinggi & 34.5$\times$ \\
45 & 16.0 & Sangat Tinggi & 80.0$\times$ \\
\bottomrule
\end{tabular}
\end{table}

Peningkatan risiko yang drastis ini selaras dengan temuan Hook (1981) dan
Morris \textit{et al.} (2002), mengkonfirmasi validitas model yang digunakan.

\subsection{Distribusi Non-Disjunction: Meiosis I vs. Meiosis II}

Hasil simulasi pada 10.000 iterasi usia 35 tahun menunjukkan distribusi:

\begin{enumerate}[leftmargin=*,noitemsep,topsep=2pt]
    \item \textbf{Non-disjunction Meiosis I}: $\approx 75.4\%$ dari total kejadian ND.
    \item \textbf{Non-disjunction Meiosis II}: $\approx 24.6\%$ dari total kejadian ND.
\end{enumerate}

Proporsi ini konsisten dengan literatur (Hassold \& Hunt 2001) yang melaporkan
dominasi Meiosis I sebesar $75$--$80\%$. Secara biologis, hal ini dijelaskan oleh
penuaan protein kohesin pada kromosom oosit: ikatan kohesin yang menjaga pasangan
kromosom homolog bersama mulai mengalami degradasi setelah puluhan tahun dalam
arrest meiosis, sehingga pemisahan prematur lebih sering terjadi di Meiosis I.

\subsection{Distribusi Sindrom Klinis}

Dari total gamet aneuploid yang dihasilkan pada usia 38 tahun ($N=50.000$),
distribusi prediksi sindrom adalah sebagai berikut:

\begin{enumerate}[leftmargin=*,noitemsep,topsep=2pt]
    \item \textbf{Trisomi 21 (Sindrom Down)}: $\approx 62.6\%$ dari kasus aneuploid,
        konsisten dengan insidensi klinis terbanyak.
    \item \textbf{Trisomi 18 (Sindrom Edwards)}: $\approx 20.0\%$
    \item \textbf{Trisomi 13 (Sindrom Patau)}: $\approx 10.0\%$
    \item \textbf{Aneuploidi Kromosom Seks}: $\approx 7.4\%$
\end{enumerate}

Distribusi ini mencerminkan bobot probabilitas pada kromosom yang lebih kecil
(kromosom 21 memiliki jumlah gen terkecil di antara autosom, sehingga lebih
``toleran'' terhadap kelebihan salinan --- gamet dengan Trisomi 21 lebih sering
menghasilkan janin viabel).

\subsection{Visualisasi 3D sebagai Kontribusi Edukatif}

Platform visualisasi Three.js berhasil merender animasi proses meiosis secara
interaktif, meliputi transisi fase dari interfase hingga pembentukan gamet.
Kegagalan non-disjunction divisualisasikan dengan kromosom berwarna merah berdenyut
yang gagal bergerak ke kutub yang berlawanan selama anafase. Gamet aneuploid
direpresentasikan sebagai bola merah dalam kumpulan gamet hasil simulasi,
dengan proporsi visual yang secara kuantitatif mencerminkan nilai risiko terobservasi.

% ══════════════════════════════════════════════════════════════════
\section{Kesimpulan}

Penelitian ini berhasil membangun sistem komputasi terpadu yang menggabungkan
pemodelan probabilistik, simulasi Monte Carlo, dan visualisasi 3D interaktif
untuk memprediksi dan memvisualisasikan aneuploidi kromosom akibat kegagalan
non-disjunction pada proses meiosis.

Kontribusi utama penelitian ini secara objektif adalah:

\begin{enumerate}[leftmargin=*,noitemsep,topsep=2pt]
    \item \textbf{Model probabilistik berbasis empiris}: model yang dikembangkan
        berkonvergensi terhadap data klinis terpublikasi dengan akurasi $>92\%$,
        mengkonfirmasi bahwa simulasi Monte Carlo efektif sebagai pendekatan
        komputasional dalam biologi.
    \item \textbf{Reprodusibilitas ilmiah}: penggunaan \textit{random seed} tetap
        memastikan hasil simulasi sepenuhnya reprodusibel, memenuhi standar
        keterbukaan penelitian (\textit{open science}).
    \item \textbf{Platform edukatif}: visualisasi Three.js merupakan kontribusi
        baru dalam pendidikan genetika --- menjembatani konsep biologi abstrak
        (segregasi kromosom) dengan representasi visual yang intuitif dan interaktif.
    \item \textbf{Dataset sintetis terbuka}: dataset 1.000 rekaman prenatal yang
        dihasilkan dapat menjadi \textit{baseline} bagi penelitian lebih lanjut.
\end{enumerate}

Beberapa saran untuk penelitian di masa depan:

\begin{enumerate}[leftmargin=*,noitemsep,topsep=2pt]
    \item Mengintegrasikan faktor risiko tambahan (indeks massa tubuh, paritas,
        riwayat kelainan kromosom keluarga, paparan lingkungan) ke dalam model
        probabilistik multivariat.
    \item Memvalidasi model menggunakan dataset \textit{chorionic villus sampling}
        (CVS) atau amniosentesis dari rekam medis nyata untuk meningkatkan akurasi klinis.
    \item Mengembangkan model \textit{agent-based} yang mensimulasikan dinamika
        molekuler kompleks kohesin dan \textit{spindle assembly checkpoint} (SAC)
        secara mekanik.
    \item Memperluas cakupan ke analisis \textit{paternal age effect} pada
        non-disjunction spermatogenik yang saat ini belum dimodelkan.
\end{enumerate}

% ══════════════════════════════════════════════════════════════════
\begin{thebibliography}{9}
\scriptsize
\bibitem{alberts2015}
Alberts, B., Johnson, A., Lewis, J., Morgan, D., Raff, M., Roberts, K., \& Walter, P. (2015).
\textit{Molecular Biology of the Cell} (6th ed.). Garland Science.

\bibitem{hook1981}
Hook, E. B. (1981). Rates of chromosome abnormalities at different maternal ages.
\textit{Obstetrics \& Gynecology}, 58(3), 282--285.

\bibitem{morris2002}
Morris, J. K., Mutton, D. E., \& Alberman, E. (2002). Revised estimates of the maternal age
specific live birth prevalence of Down's syndrome.
\textit{Journal of Medical Screening}, 9(1), 2--6.

\bibitem{hassold2001}
Hassold, T., \& Hunt, P. (2001). To err (meiotically) is human: the genesis of human
aneuploidy. \textit{Nature Reviews Genetics}, 2(4), 280--291.

\bibitem{hunt2008}
Hunt, P. A., \& Hassold, T. J. (2008). Human female meiosis: what makes a good egg go bad?
\textit{Trends in Genetics}, 24(2), 86--93.

\bibitem{metropolis1949}
Metropolis, N., \& Ulam, S. (1949). The Monte Carlo method.
\textit{Journal of the American Statistical Association}, 44(247), 335--341.

\bibitem{acmg2016}
ACMG. (2016). \textit{Technical standards and guidelines for reproductive genetics}.
American College of Medical Genetics and Genomics.

\bibitem{savva2010}
Savva, G. M., Walker, K., \& Morris, J. K. (2010). The maternal age-specific live birth prevalence of trisomies 13 and 18 compared to trisomy 21 (Down syndrome). \textit{Prenatal Diagnosis}, 30(1), 57--64.

\bibitem{unsw}
UNSW Embryology. \textit{Genetic Risk Maternal Age Table}.
Diakses dari: \url{https://embryology.med.unsw.edu.au/embryology/index.php?title=Template:Genetic_risk_maternal_age_table}

\bibitem{omim}
Online Mendelian Inheritance in Man (OMIM). \textit{Katalog Penyakit Genetik (Trisomi 21, 18, 13, Turner, Klinefelter)}.
Diakses dari: \url{https://www.omim.org/entry/190685}
\end{thebibliography}

\end{document}
